\section{The Mechanism of Trust}
\label{p13:sec:mechanism}

\noindent
Everything to this point has cleared the ground for a single question and supplied the terms in which it can be asked. The symptomatology displayed a gap between the apparatus of assurance and the trust it failed to secure; the diagnosis traced the gap to a mistake about causation, the taking of the mechanical cause for the cause of an inferential effect; the jurisprudence showed the law's own anchor migrating, across a century, from force toward trust; and the preceding section established that the inferential cause is a different kind of thing from the mechanical, the kind that addresses a free being as free. What none of these has yet supplied is the positive account: \emph{by what mechanism} a free mind, addressed as free and given the right grounds, comes to expect that a promise about a non-existent future, carrying no guarantee of consequence, will be kept. This section supplies it. It is the trunk of the paper, of which the chapters before were the roots and the chapters after are the branches.

The problem the mechanism must solve has, recall, a shape that looks impossible. Trust is directed at a future that does not yet exist and cannot be verified; the symbol that asks for trust carries no force and so cannot secure the future by threatening a consequence. We are asked to explain how a present expectation, real and often well-founded, is produced by a symbol about what does not exist and under no guarantee\,---\,a generation, very nearly, from nothing. The dissolution of the apparent impossibility is the section's first move, and it is made by reaching for the one vocabulary in the paper's possession that is native to the forming of expectations about what has not yet arrived: the predictive account of the mind.

\subsection{Predictive coding and the dissolution of the paradox}
\label{p13:sub:mech-pc}

The brain, on the predictive account that has become central to the cognitive sciences, is not in the first instance a device that registers the world and then responds to it; it is a device that \emph{predicts} the world and corrects its predictions against what arrives~\parencite{friston2010,clark2013,hohwy2013}. It carries a generative model of the causes of its sensory states, continually generates expectations of what is to come, compares them against incoming signals, and updates the model in proportion to the error\,---\,and in proportion, crucially, to the \emph{precision}, the inverse variance, it assigns to that error and to its own predictions. Prediction, on this account, is not an occasional achievement but the standing activity of the organ; the mind is forward-directed by its nature, expecting before it perceives.

This dissolves the apparent paradox at a stroke. The objection was that trust is impossible because its object\,---\,the kept future\,---\,does not yet exist; how can a present state be produced by an absent fact? But the objection assumes that an expectation must be caused by the thing expected, and the predictive account shows this assumption false. An expectation is not caused by its object; it is produced, \emph{now}, by the generative model, as the model's forward prediction, in advance of and independent of the object's arrival. Forward-directed expectation under uncertainty is not a special feat trust must somehow accomplish against the grain of cognition; it is what cognition is incessantly doing anyway. Trust requires no time machine, no causation running backward from a future that does not exist. It requires only what the brain does in every waking moment: the assignment, by a generative model, of a high-precision prediction\,---\,here, the prediction that this person will keep this promise. \emc{Trust is the high precision a receiver's generative model assigns to the prediction that the promise will be kept.} The expectation of legal possibility that the diagnosis spoke of is, at the level of mechanism, exactly this: a precision-weighted forward prediction, a quantity the brain natively computes, and not a metaphor borrowed from the law.

The question of how trust is generated thus becomes a precise question: by what \emph{inputs} does a receiver's generative model come to assign high precision to the prediction that a promise will be kept? The paper's answer is that there are three such inputs, three channels, and that they correspond\,---\,not by coincidence but because the registers are the registers of one subject\,---\,to the three registers the series has worked with throughout.

\subsection{The three channels}
\label{p13:sub:mech-channels}

The three channels are distinct in what they supply to the generative model, and the paper takes them in the order of the registers: the structural, lodged in the symbolic; the self-committing, lodged in the real; and the wagered, lodged in the imaginary and the subject who wagers.

\subsubsection{The structural channel: the prior}
\label{p13:sub:mech-prior}

The first channel supplies a \emph{prior}. A promise is not received in a vacuum; it is received with a structure, and the structure of the promise shapes the receiver's prior expectation before any evidence about this particular promisor is weighed. A vow with the form of a closed and self-completing commitment\,---\,one that, in the language of \xref{p13:sec:holonomy}, returns upon itself without leaving a debt to be discharged elsewhere\,---\,presents to the receiver's model a structural prior favourable to its being kept, in a way that a loose, hedged, escape-laden form does not. This is the cognitive function of what the series has called the poetic structure of a commitment: the form is not ornament but information, a shaper of the prior the receiver brings to bear. At the scale of institutions this is the channel of Hart's convergent acceptance and the rule of recognition (\xref{p13:sub:jur-hart}): a shared structural prior, sustained by being shared, against which each new undertaking is read. The structural channel is lodged in the symbolic, for it works through the form of the symbol; and it is, taken alone, the most forgeable of the three, since a form can be imitated.

\subsubsection{The self-committing channel: the high-fidelity likelihood}
\label{p13:sub:mech-likelihood}

The second channel supplies a \emph{likelihood}\,---\,and, where it works, a high-fidelity one. Here the promisor does not merely present a form; he does something. He places himself, really and observably, in a position of vulnerability from which retreat is costly or foreclosed: he binds himself in a manner that will hurt him to break, gives a hostage, burns a bridge, commits himself irrevocably and in the sight of others. And the receiver's generative model reads this for what, statistically, it is. A signal that is costly and hard to reverse carries a high precision-weight, because it is one a promisor who did \emph{not} intend to keep faith would be very unlikely to send: the cost of sending it, for the insincere, is prohibitive, and so its being sent is strong evidence of sincerity. This is the mechanism beneath the red silk, the public setting, the burning of one's boats\,---\,the irrevocable and costly self-binding the Prelude began with. Such acts are not arguments; they do not persuade by inference from premises. They furnish the receiver's model with a high-fidelity, hard-to-forge likelihood signal, evidence of a kind the body reads before the discourse does, that the one who could so bind himself must mean it.

The paper takes up here, and gives a mechanistic basis to, the costly-signalling principle familiar from the theory of signalling in biology and economics: that a signal too expensive for a defector to afford is, for that reason, credible. What predictive coding adds is the implementation\,---\,the costly, irrevocable signal is precisely one to which a generative model rationally assigns high precision\,---\,and a deeper observation the paper will not exhaust here but must mark. The burning of the bridge has been described, just now, in its aspect as a \emph{signal}: as something done to be read by another, furnishing his model with evidence. But this is only one of its two faces, and not the deeper. There is a second sense in which one burns a bridge that has nothing to do with being read by anyone\,---\,in which one forecloses one's own retreat not to send a message but to constitute, for oneself, a law: \emph{I have decreed it so, and so that road no longer exists for me}. This deeper sense, in which the irrevocable self-binding is an act of self-legislation rather than of signalling, is the matter of \xref{p13:sec:praxis}, and in particular of its account of relational divinity; it is marked here, at the mechanism, so that the reader who reaches it will recognise that the trunk was pointing toward it all along. For the present chapter the self-committing channel is treated in its signalling aspect: as the high-fidelity likelihood lodged in the real, the channel through which a real act of self-exposure, received by another present to read it, supplies the model with its least forgeable evidence.

\subsubsection{The wagered channel: the prior leap of active inference}
\label{p13:sub:mech-wager}

The third channel is the most easily mistaken for irrationality and is, properly understood, the deepest. It concerns not what the promisor supplies but what the receiver does: the receiver, in trusting, makes a leap that precedes the evidence. He assigns, in advance of sufficient warrant, a high-precision prior to the prediction that the other will keep faith\,---\,he ventures the trust\,---\,and only then, through the continued interaction the venture makes possible, gathers the error signals by which the model is confirmed or corrected. This looks like a failure of rationality: the trust is given before the evidence that would justify it. But on the active-inference reading it is not a failure of rationality at all. An agent does not merely passively update a model against incoming data; it acts so as to reduce its own uncertainty over the longer run, and there are states of knowledge attainable only by first acting under uncertainty\,---\,only by venturing. Trust is such a venture. One cannot gather the evidence that would warrant trusting a person without first entering the kind of interaction that only trusting him makes possible; the evidence lies on the far side of a threshold one can cross only by trusting in advance of it. The leap of trust is thus the rational opening move of an agent that reduces uncertainty by acting\,---\,a prior wager that is the precondition of acquiring the very evidence that will vindicate or correct it. This channel is lodged in the imaginary and in the wagering subject, for its source is not in the symbol's form nor in the promisor's act but in the receiver's free venturing; and it is here, in the wager that the other will prove free and faithful, that the freedom of \xref{p13:sec:freedom} re-enters the mechanism.

\subsection{The multiplicative coupling}
\label{p13:sub:mech-coupling}

The three channels are not three alternative routes to trust, of which a promisor might supply one and be believed. They are coupled, and coupled \emph{multiplicatively}\,---\,as the factors of a single Bayesian product are coupled, in which the structural prior, the self-committing likelihood, and the wagered prior combine into the posterior precision of the prediction that the promise will be kept. The point of multiplicative rather than additive coupling is exact and it is the heart of the section: in a product, if any factor goes to zero the whole product goes to zero, and no abundance in the other factors can compensate. A perfect form with no real self-commitment behind it; a costly act with no wager that the other is free to be faithful; a leap of trust met by neither form nor self-binding\,---\,each is a product with a vanishing factor, and each fails as a ground of well-founded trust.

This is why genuine trust is robust where any single channel would be fragile. A single channel can be forged: a form can be imitated, a costly-looking act can be staged, a credulous wager can be exploited. But the three channels are hard to forge \emph{together}, because they are answerable to one another\,---\,the structural form must be borne out by a real self-commitment, which must in turn be the kind of thing a free agent would do and would invite a wager upon\,---\,and their mutual consistency, the agreement of form with act with wager, is far harder to counterfeit than any one of them alone. The strength of genuine trust is not the strength of its strongest channel but the improbability of all three being made to agree falsely. This is what coupling means here, and why it is the source of trust that is not merely felt but well-founded.

\begin{claim}[The mechanism of trust]
The expectation that a non-coercive promise will be kept is the posterior precision a receiver's generative model assigns to that prediction, formed by the multiplicative coupling of three inputs: a structural prior, supplied by the form of the commitment (the symbolic register); a high-fidelity likelihood, supplied by the promisor's real, costly, irrevocable self-binding (the real register); and a prior wager, supplied by the receiver's free venturing of trust in advance of the evidence (the imaginary register, and the wagering subject). Genuine trust is the state in which all three are present, real, and mutually consistent; and its well-foundedness rests on the improbability of all three being made to agree falsely.
\end{claim}

\subsection{The forged case, mechanically described}
\label{p13:sub:mech-forged}

The mechanism furnishes, as a by-product, an exact characterisation of its own counterfeit, which \xref{p13:sec:justice} will take up as the matter of justice and which is set out here in its mechanical form. Forged trust is the state in which the receiver's generative model is brought to assign high precision to the prediction that the promise will be kept \emph{without} the real self-committing likelihood that alone would warrant it. The manipulator does precisely this: by working on the receiver's model directly\,---\,by affective pressure, by the staging of forms, by rhetorical management\,---\,he induces a lowering of the receiver's prediction error, a high-precision prediction, that is not supported by any real self-binding on his part. The precision is, as it were, injected from outside rather than earned by a genuine likelihood; and a precision so injected is a prediction unsupported by what it predicts, a confidence with nothing real beneath it.

The mechanical mark of manipulation can therefore be stated with some sharpness: \emb{to manipulate is to induce the other's generative model to lower its prediction error without supplying the high-fidelity likelihood that would warrant the lowering.} Such a state is unstable in a way genuine trust is not. Because no real self-commitment underlies it, the world will not in fact behave as the injected precision predicts; the error signals will accumulate, the model will be forced toward correction, and the divergence between the predicted and the actual will widen until the forged precision collapses\,---\,the process the series has elsewhere described, in its dynamical register, as the accumulation of deviation toward divergence under a non-adiabatic extraction. The forged case is thus not a stable alternative form of trust but a transient, a precision held up by pressure rather than by likelihood, destined for the collapse that the absence of a real ground guarantees.

\subsection{The crisis of the symbol, in predictive-coding terms}
\label{p13:sub:mech-crisis}

The same vocabulary yields, finally, a precise rendering of the diagnosis of \xref{p13:sec:diagnosis}, and shows the crisis of the symbol to be not a moral decline but a rational adjustment. Consider the population of generative models that is a society, each learning, from its own commerce with symbols, the statistical relation between symbolic signals and real performance. As the cost of producing a convincing symbol falls toward zero and forged signals proliferate (the generativity paradox of \xref{p13:sub:diag-three}), each model learns, correctly, that the correlation between symbolic testimony and real performance has weakened\,---\,that the signal has become, in fact, less predictive of the outcome. And the rational Bayesian response to a signal that has become less predictive is to \emph{lower the precision-weight assigned to it}. Each model, updating as it should, down-weights symbolic testimony as such; and because the down-weighting is applied to the channel and not to the individual signal, it falls upon the genuine symbol along with the forged. The genuine vow, the true recommendation, the sincere assurance\,---\,each is met by a model that has rationally learned to trust symbols of its kind less, and so each fails to land, not because it is disbelieved in particular but because the channel it travels has been rationally discounted.

This is the crisis of the symbol given its mechanism. It is not that people have become cynical, or that values have decayed; it is that a population of well-functioning generative models, faced with a channel that forgery has made less reliable, has rationally collapsed the precision-weight it assigns to that channel\,---\,and that the retreat to the merely countable (\xref{p13:sub:sympt-universal}) is the same population's attempt to reconstruct a high-fidelity likelihood out of signals that are, at least, harder to forge. The crisis is a rational precision collapse at the scale of a population, and that is both why it is so intractable\,---\,for each model is only doing what it should\,---\,and why the remedy, as the rest of the paper develops it, cannot be a return to symbolic profession but must lie in the channels, the self-committing and the relational, that forgery cannot so easily saturate.

\subsection{Connection to the computational programme of the foundational paper}
\label{p13:sub:mech-Jbracket}

It is worth marking, before the chapter closes, that the mechanism dovetails with the formal programme the foundational paper held open. That paper sketched, as a direction for future work, a computational study of justice as generative presence, in which justice would figure as an observable evaluated over a space of histories, formalised in the manner of a path integral~\parencite{wanhong_grb}. The predictive-coding mechanism of trust is a concrete realisation, for the case of trust, of exactly this form: trust is the precision a generative model assigns to a prediction evaluated over the space of the promisor's possible future trajectories, an expectation taken over histories. The two formal programmes mesh\,---\,the observable-over-histories of the foundational paper and the precision-over-trajectories of the present mechanism are the same kind of object\,---\,and the meshing is noted here so that the formal continuity of the series is on the record, to be drawn upon rather than re-derived.

\subsection{A discipline of honesty}
\label{p13:sub:mech-honesty}

The chapter must end where the series requires it to, with an express discipline upon the use it has made of the predictive account, lest the mechanism be taken for more than it is. Predictive coding has been used here as a \emph{mechanistic and structural model}, not as a reductive claim about what trust ultimately is. The paper does not assert that trust simply \emph{is} a precision-weighting and nothing more. What predictive coding supplies is a description at Marr's computational and implementational levels\,---\,an account of what the system computes and how it might compute it\,---\,and such a description does not displace the others the paper has given and will give: the phenomenological description of the vulnerability of self-commitment, the jurisprudential description of force as a substitute for expectation, the ethical description of the difference between genuine and forged. These are descriptions of one object at different levels, and the predictive-coding account is one of them, not the truth beneath them all. It is not a meta-framework that absorbs the rest; it is one more irreducible position on a single object, and the polyphonic discipline of the series\,---\,that no framework govern the others\,---\,holds over it as over every other. The mechanism is offered, then, as the paper's central contribution and not as its last word: a specification of how, at the level of the computing mind, the inferential cause does its work, which leaves standing, beside it, everything the other registers have to say about the same thing.
